% titel: Venturi-Rohr
% skript: Kap. Hydrodynamik

% Teil I: Werkzeuge, kein Lösungsweg.

% Ein Venturi-Messrohr wandelt eine Geschwindigkeitsdifferenz (wegen verengtem
% Querschnitt) in eine Druckdifferenz um; gesucht ist der Volumenstrom aus dem
% gemessenen Druckunterschied (hier über ein Quecksilbermanometer).

Ein reibungsfrei, inkompressibel strömendes Fluid passiert eine Verengung; die
Druckdifferenz zwischen zwei Querschnitten wird gemessen (hier über ein
U-Rohr-Manometer mit Quecksilber). Gesucht ist der Volumenstrom als Funktion
der Durchmesser und der Manometer-Höhendifferenz.

\begin{werkzeugkasten}[Werkzeuge]
\textbf{Kontinuitätsgleichung} (Masseerhalt, inkompressibel):
\begin{equation}
  A_1 v_1 = A_2 v_2 = \dot V \;,\qquad A=\frac{\pi d^2}{4}\;.
\end{equation}

\textbf{Bernoulli-Gleichung} (Energieerhaltung, stationär, reibungsfrei):
\begin{equation}
  p_1+\frac{\rho}{2}\,v_1^2+\rho g z_1
  =\; p_2+\frac{\rho}{2}\,v_2^2+\rho g z_2\;.
\end{equation}
Bei horizontalem Rohr ($z_1=z_2$) bleibt nur die Umwandlung
Druckenergie $\leftrightarrow$ kinetische Energie:
\begin{equation}
  p_1-p_2=\frac{\rho}{2}\left(v_2^2-v_1^2\right)\;.
\end{equation}

\textbf{U-Rohr-Manometer:} Die Druckdifferenz wird über eine Säulenhöhe
$\Delta h$ der Manometer-Flüssigkeit (Dichte $\rho_\mathrm{m}$) abgelesen,
\begin{equation}
  p_1-p_2=\rho_\mathrm{m}\,g\,\Delta h\;.
\end{equation}

\textbf{Einheiten und Bezeichner:} $\dot V$ in \unit{\cubic\metre\per\second},
Durchmesser $d$ in \unit{\metre}, $\rho$ in \unit{\kilogram\per\cubic\metre},
$g\approx\qty{9.81}{\metre\per\square\second}$.
\end{werkzeugkasten}

\begin{werkzeugkasten}[Vorgehen in allgemeiner Form]
Aus Bernoulli folgt der Geschwindigkeitsunterschied aus der Druckdifferenz,
aus der Kontinuität der Zusammenhang der beiden Querschnitts-Geschwindigkeiten.
Kombiniert man beide und löst nach einer Geschwindigkeit auf, hängt diese nur
von Querschnittsverhältnis und Manometerhöhe ab. Der Volumenstrom ist
Geschwindigkeit mal Fläche.
\end{werkzeugkasten}

\begin{achtung}[Typische Fallen]
\begin{itemize}
  \item Zwei Dichten nicht verwechseln: $\rho$ des \emph{strömenden Gases} gegen
       $\rho_\mathrm{m}$ des \emph{Manometer-Quecksilbers}.
  \item Die Druckdifferenz des Manometers ist $p_1-p_2$, nicht $p_2-p_1$.
  \item Im Volumenstrom-Ausdruck erscheint $d_1$ und $d_2$ als Quotient,
       also über $A_2/A_1=(d_2/d_1)^2$.
\end{itemize}
\end{achtung}
